% ===================================================================
%  ANTENNA PATTERN IMPLEMENTATION ANALYSIS
%  Comparative Study: AntennaForge (Python) vs ODESSA (C++)
%
%  Compile with:  pdflatex antenna_implementation_report.tex
% ===================================================================
\documentclass[11pt,a4paper]{article}

% ── Packages ──────────────────────────────────────────────────────
\usepackage[utf8]{inputenc}
\usepackage[T1]{fontenc}
\usepackage{lmodern}
\usepackage[margin=1in]{geometry}
\usepackage{amsmath,amssymb,amsfonts}
\usepackage{booktabs}
\usepackage{array}
\usepackage{longtable}
\usepackage{enumitem}
\usepackage{xcolor}
\usepackage{hyperref}
\usepackage{titlesec}
\usepackage{fancyhdr}
\usepackage{graphicx}
\usepackage{float}
\usepackage{caption}
\usepackage{subcaption}
\usepackage{listings}
\usepackage{multirow}
\usepackage{tabularx}

% ── Colour definitions ────────────────────────────────────────────
\definecolor{codegreen}{rgb}{0.0, 0.5, 0.0}
\definecolor{codegray}{rgb}{0.5, 0.5, 0.5}
\definecolor{codepurple}{rgb}{0.58, 0.0, 0.82}
\definecolor{backcolour}{rgb}{0.97, 0.97, 0.97}
\definecolor{sectionblue}{rgb}{0.1, 0.2, 0.5}
\definecolor{highlightgreen}{rgb}{0.85, 0.95, 0.85}
\definecolor{highlightred}{rgb}{0.95, 0.85, 0.85}
\definecolor{highlightyellow}{rgb}{0.95, 0.95, 0.80}

% ── Code listing style ────────────────────────────────────────────
\lstdefinestyle{codestyle}{
    backgroundcolor=\color{backcolour},
    commentstyle=\color{codegreen}\itshape,
    keywordstyle=\color{codepurple}\bfseries,
    numberstyle=\tiny\color{codegray},
    stringstyle=\color{codepurple},
    basicstyle=\ttfamily\footnotesize,
    breakatwhitespace=false,
    breaklines=true,
    captionpos=b,
    keepspaces=true,
    numbers=left,
    numbersep=5pt,
    showspaces=false,
    showstringspaces=false,
    showtabs=false,
    tabsize=2,
    frame=single,
    rulecolor=\color{codegray},
}
\lstset{style=codestyle}

% ── Section formatting ────────────────────────────────────────────
\titleformat{\section}
  {\Large\bfseries\color{sectionblue}}{\thesection}{1em}{}
\titleformat{\subsection}
  {\large\bfseries\color{sectionblue!80}}{\thesubsection}{1em}{}
\titleformat{\subsubsection}
  {\normalsize\bfseries\color{sectionblue!60}}{\thesubsubsection}{1em}{}

% ── Header / footer ──────────────────────────────────────────────
\pagestyle{fancy}
\fancyhf{}
\fancyhead[L]{\small\textit{Antenna Pattern Implementation Analysis}}
\fancyhead[R]{\small\thepage}
\fancyfoot[C]{\small\textit{AntennaForge vs ODESSA Comparative Study}}
\renewcommand{\headrulewidth}{0.4pt}
\renewcommand{\footrulewidth}{0.4pt}

% ── Custom commands ───────────────────────────────────────────────
\newcommand{\AF}{\textsc{AntennaForge}}
\newcommand{\OD}{\textsc{ODESSA}}
\newcommand{\sinc}{\operatorname{sinc}}
\newcommand{\caveat}[1]{\par\noindent\colorbox{highlightyellow}{\parbox{\dimexpr\linewidth-2\fboxsep\relax}{\textbf{Note:} #1}}\par\medskip}

% ── Hyperref setup ────────────────────────────────────────────────
\hypersetup{
    colorlinks=true,
    linkcolor=sectionblue,
    citecolor=sectionblue,
    urlcolor=sectionblue!70,
    pdftitle={Antenna Pattern Implementation Analysis},
    pdfauthor={AntennaForge},
}

% ══════════════════════════════════════════════════════════════════
\begin{document}

% ── Title page ────────────────────────────────────────────────────
\begin{titlepage}
\centering
\vspace*{3cm}
{\Huge\bfseries\color{sectionblue} Antenna Pattern\\[0.3em]
Implementation Analysis\par}
\vspace{1.5cm}
{\Large Comparative Study:\\[0.3em]
\AF{} (Python) vs \OD{} (C++)\par}
\vspace{2cm}
{\large February 2026\par}
\vfill
{\small
This document provides a detailed mathematical comparison of antenna
radiation pattern models implemented in the \AF{} Python framework
and the \OD{} C++ simulation library. Each antenna type is examined
for correctness, physical fidelity, numerical robustness, and
feature completeness.\par}
\vspace{1cm}
\end{titlepage}

% ── Table of contents ─────────────────────────────────────────────
\tableofcontents
\newpage

% ══════════════════════════════════════════════════════════════════
\section{Executive Summary}
% ══════════════════════════════════════════════════════════════════

This report analyses seven antenna types implemented across two independent codebases:

\begin{itemize}[leftmargin=2em]
  \item \textbf{\AF{}} --- a Python-based antenna pattern generation framework using
        Jinc/Airy, pyramidal horn aperture theory, and full array factor
        summation models, with extensive post-processing features.
  \item \textbf{\OD{}} --- a C++ antenna library designed for real-time system-level
        simulation, with multi-beam management, monopulse processing, and
        lookup-table pattern support.
\end{itemize}

\noindent The key findings are summarised in Table~\ref{tab:summary}.

\begin{table}[H]
\centering
\caption{High-level comparison of pattern models.}
\label{tab:summary}
\small
\begin{tabularx}{\textwidth}{l X X c}
\toprule
\textbf{Antenna} & \textbf{\AF{} Model} & \textbf{\OD{} Model} & \textbf{Availability} \\
\midrule
Dish (circular) & Jinc$^2$ (Airy pattern) & sinc (rect.\ aperture) & Both \\
Horn (pyramidal) & E: sinc$^2$, H: cos-taper & cos($\theta$) approx. & Both \\
Dipole / Monopole & Classical thin-wire + normalisation & Classical thin-wire & Both \\
Phased Array & Element $\times$ AF phasor sum & Parametric sinc aperture & Both \\
Panel / Sector & sinc$^2$ rectangular aperture & --- & \AF{} only \\
LPDA & Gaussian + radial sidelobes & --- & \AF{} only \\
Isotropic & (Omni cos$^n$ model) & $G = 1$ & Both \\
Elliptical Dish & --- & sinc separable $W \times H$ & \OD{} only \\
Lookup (table) & --- & 2-D / 3-D interpolation & \OD{} only \\
\bottomrule
\end{tabularx}
\end{table}

% ══════════════════════════════════════════════════════════════════
\section{Architectural Comparison}
% ══════════════════════════════════════════════════════════════════

\subsection{Design Philosophy}

\begin{description}[leftmargin=1.5em, style=nextline]
  \item[\AF{}] Prioritises \emph{pattern fidelity} and \emph{extensibility}. Each
    antenna type implements a single \\ \texttt{compute\_point\_gain()} method that
    returns the co-pol linear gain at one $(az, el)$ point. All system-level
    effects---frequency-dependent beamwidth, VSWR rolloff, asymmetry, sidelobes,
    pattern breakup, ground reflection, cross-polarisation, and noise---are
    handled by the engine before or after the per-point call.

  \item[\OD{}] Prioritises \emph{runtime speed} and \emph{system integration}.
    Each antenna implements \texttt{CalculateTxGain()} and
    \texttt{CalculateRxGain()} with monopulse ($\Sigma$, $\Delta_{az}$,
    $\Delta_{el}$) channels. Multi-beam management (add, resteer, remove) and
    spoiled transmit beams are natively supported.
\end{description}

\subsection{Coordinate Conventions}

\begin{table}[H]
\centering
\caption{Coordinate and unit conventions.}
\small
\begin{tabularx}{\textwidth}{l X X}
\toprule
\textbf{Aspect} & \textbf{\AF{}} & \textbf{\OD{}} \\
\midrule
Angle input & Degrees (converted internally) & Radians (custom unit types) \\
Boresight & $az = 0$, $el = 0$ (broadside) & $az = 0$, $el = 0$ (along axis for dipoles) \\
Gain units & Linear (engine converts to dBi) & Dimensionless linear \\
Frequency & MHz throughout & Hz (SI unit type) \\
Physical constants & $c = 300/f_\text{MHz}$ (wavelength) & \texttt{speed\_of\_light} constant \\
\bottomrule
\end{tabularx}
\end{table}

% ══════════════════════════════════════════════════════════════════
\section{Parabolic Dish Antenna}
\label{sec:dish}
% ══════════════════════════════════════════════════════════════════

\subsection{Physical Background}

A parabolic reflector with a circular aperture of diameter $D$ produces a
far-field radiation pattern governed by the Fourier transform of the circular
aperture distribution.  For uniform illumination, the exact result is the
\emph{Airy pattern}:
\begin{equation}
  G(\theta) = G_{\text{peak}} \cdot \eta \cdot
    \left[ \frac{2\,J_1(u)}{u} \right]^2,
  \qquad
  u = \pi \frac{D}{\lambda} \sin\theta
  \label{eq:airy}
\end{equation}
where $J_1$ is the Bessel function of the first kind, order one, and $\eta$ is
the aperture efficiency (typically 0.55--0.65).

The quantity $[2J_1(u)/u]^2$ is known as the \emph{Jinc-squared} or Airy
function.  Its first null occurs at $u \approx 3.832$ and first sidelobe at
$u \approx 5.136$ with an amplitude of $-17.6$~dB relative to the peak.

\subsection{\AF{} Implementation}

\AF{} implements the full Airy/Jinc$^2$ model in \texttt{dish.py}:

\begin{equation}
  G(\theta_{az}, \theta_{el}) = G_{\text{peak}} \cdot \eta \cdot
    \left[ \frac{2\,J_1(u)}{u} \right]^2 \cdot f_{\text{front}}(\cos\theta_{az}),
  \qquad
  u = K_{3\text{dB}} \cdot r
  \label{eq:af_dish}
\end{equation}
where
\[
  r = \sqrt{
    \left(\frac{\sin\theta_{az}}{\sin(\theta_{az,3\text{dB}}/2)}\right)^2 +
    \left(\frac{\sin\theta_{el}}{\sin(\theta_{el,3\text{dB}}/2)}\right)^2
  }
\]
is the normalised elliptical radial distance, and $K_{3\text{dB}} = 1.6163$ is
chosen such that $[\,2J_1(K)/K\,]^2 = 0.5$ (exactly $-3$~dB at the half-power
beamwidth).

Key implementation details:
\begin{itemize}[leftmargin=2em]
  \item \textbf{Bessel function:} Uses \texttt{scipy.special.j1} when available;
        falls back to an 8-significant-figure polynomial approximation
        (Abramowitz \& Stegun) for $|x| < 8$ and asymptotic form for larger
        arguments.
  \item \textbf{Singularity at $u = 0$:} Explicitly returns $1.0$ when
        $|u| < 10^{-10}$, avoiding the $0/0$ indeterminate form.
  \item \textbf{Elliptical beam:} Supports independent azimuth and elevation
        beamwidths via the $r$ normalisation, enabling elliptical aperture
        modelling.
  \item \textbf{Front-to-back ratio:} Models rear-hemisphere radiation with a
        configurable FTB parameter and continuous $(1 - \cos\theta_{az})/2$
        back-lobe weighting.
  \item \textbf{Front fade:} Uses a sigmoid function
        $f_{\text{front}} = 1/(1 + e^{-k\cos\theta_{az}})$ with a
        user-configurable steepness $k$ (default~$40$, range~$1$--$100$)
        to smoothly transition between front and back hemispheres, eliminating
        the $\pm 90°$ discontinuity.  Lower $k$ values produce a wider,
        more physically realistic transition zone.
\end{itemize}

\subsection{\OD{} Implementation}

\OD{} models the parabolic dish in \texttt{antenna\_parabolic\_dish.cpp} using
a sinc-based pattern:
\begin{equation}
  G(\theta) = G_{\text{peak}} \cdot \eta \cdot
    \left| \sinc\!\left(\theta \cdot \pi \frac{D}{\lambda}\right) \right|,
  \qquad
  G_{\text{peak}} = \eta \left(\frac{\pi D}{\lambda}\right)^2
  \label{eq:od_dish}
\end{equation}

The off-axis angle is computed as:
\[
  \theta = \sqrt{\theta_{az}^2 + \theta_{el}^2}
\]

For the combined transmit gain, the azimuth and elevation angles off boresight
are combined into a single radial angle. Monopulse receive processing computes
difference channels:
\[
  \Delta_{az} = \sqrt{3}\,\frac{\sin(x_{az}) - x_{az}\cos(x_{az})}{x_{az}^2}
  \cdot G_{\text{peak}},
  \qquad
  x = \pi \frac{D f}{c} \cdot \theta_{\text{sat}}
\]

\subsection{Mathematical Analysis}

\subsubsection{Diffraction Kernel}

The fundamental difference is the choice of diffraction kernel:

\begin{table}[H]
\centering
\caption{Diffraction kernel comparison for circular aperture.}
\small
\begin{tabularx}{\textwidth}{l X X}
\toprule
\textbf{Property} & \textbf{Jinc$^2$ (\AF{})} & \textbf{sinc (\OD{})} \\
\midrule
Aperture geometry & Circular (correct for dish) & Rectangular (1-D projection) \\
First null location & $u = 3.832$ ($1.22\lambda/D$) & $u = \pi$ ($\lambda/D$) \\
First sidelobe & $-17.6$~dB & $-13.3$~dB \\
Second sidelobe & $-23.8$~dB & $-17.9$~dB \\
Sidelobe falloff & $\sim u^{-3}$ & $\sim u^{-2}$ \\
\bottomrule
\end{tabularx}
\end{table}

The sinc function is the Fourier transform of a \emph{rectangular} aperture.
Applying it to a circular dish produces sidelobes that are 4.3~dB too high at
the first sidelobe and that decay too slowly ($u^{-2}$ vs $u^{-3}$), resulting
in an over-estimation of interference power at wide angles.

\subsubsection{Beamwidth Scaling}

Both implementations use the standard $\theta_{3\text{dB}} = 1.22\lambda/D$
beamwidth relation for a circular aperture.  \AF{} additionally supports
frequency-dependent beamwidth scaling via three configurable modes (exponent,
percentage per octave, three-point Lagrange interpolation).

\subsubsection{Peak Gain}

Both use the standard formula:
\[
  G_{\text{peak}} = \eta \left(\frac{\pi D}{\lambda}\right)^2
\]
\OD{} computes this explicitly in \texttt{GetMaxGain()}.  \AF{} accepts the
user-specified peak gain in dBi and applies the efficiency factor as a
multiplicative weight in the pattern computation.

\caveat{The Jinc$^2$ kernel corresponds to the Fourier transform of a
circular aperture and produces first sidelobes at $-17.6$~dB.  The sinc
kernel corresponds to a rectangular aperture, with first sidelobes at
$-13.3$~dB.  This 4.3~dB difference in sidelobe level is the primary
distinction between the two dish models.}

% ──────────────────────────────────────────────────────────────────
\subsubsection{\OD{} Monopulse Receive (Feature Absent in \AF{})}

\OD{} implements monopulse difference-channel patterns for tracking radars:
\begin{align}
  \Sigma &= G_{\text{peak}} \cdot |\sinc(x)| \\
  \Delta &= \sqrt{3} \cdot G_{\text{peak}} \cdot
    \frac{\sin(x) - x\cos(x)}{x^2}
\end{align}

This is a standard approximation for a monopulse comparator on a uniformly
illuminated circular aperture.  \AF{} does not currently model monopulse
processing.

% ══════════════════════════════════════════════════════════════════
\section{Horn Antenna}
\label{sec:horn}
% ══════════════════════════════════════════════════════════════════

\subsection{Physical Background}

A pyramidal horn antenna has distinct radiation patterns in the E-plane
(elevation) and H-plane (azimuth):

\begin{align}
  \text{E-plane (uniform):} \qquad
    F_E(\theta) &= \sinc\!\left(\frac{\pi a}{\lambda}\sin\theta\right)
    \label{eq:horn_e} \\[6pt]
  \text{H-plane (cosine taper):} \qquad
    F_H(\theta) &= \frac{\cos\!\left(\frac{\pi b}{\lambda}\sin\theta\right)}
         {1 - \left(\frac{2b}{\lambda}\sin\theta\right)^2}
    \label{eq:horn_h}
\end{align}

where $a$ and $b$ are the E-plane and H-plane aperture dimensions, respectively.

For a conical horn (circular aperture), the pattern is symmetric and governed
by the Airy function---similar to a dish but with a different feed mechanism.

\subsection{\AF{} Implementation}

\AF{} implements the full \emph{pyramidal} horn model with separate E-plane
and H-plane functions:

\begin{equation}
  G(\theta_{az}, \theta_{el}) = G_{\text{peak}} \cdot
    F_H^2(v_{az}) \cdot F_E^2(u_{el}) \cdot f_{\text{front}}
\end{equation}
where
\begin{align}
  u_{el} &= K_E \cdot \frac{\sin\theta_{el}}{\sin(\theta_{el,3\text{dB}}/2)},
  &\qquad K_E &= 1.3916 \quad (\sinc^2(K_E) = 0.5) \\
  v_{az} &= K_H \cdot \frac{\sin\theta_{az}}{\sin(\theta_{az,3\text{dB}}/2)},
  &\qquad K_H &= 1.1891
\end{align}

The H-plane function includes proper treatment of the singularity at
$v = \pm\pi/2$ via L'H\^{o}pital's rule:
\[
  \lim_{v \to \pm\pi/2} F_H(v) = \frac{\pi}{4}\cos(v)
\]

\subsection{\OD{} Implementation}

\OD{} models a \emph{conical} horn in \texttt{antenna\_horn.cpp} using a
simple cosine approximation:
\begin{equation}
  G(\theta) = G_{\text{peak}} \cdot \sin\!\left(\frac{\pi}{2} - \theta\right)
            = G_{\text{peak}} \cdot \cos(\theta)
\end{equation}
where $\theta$ is computed as the radial angle off boresight:
\[
  \theta = \sqrt{\theta_{az}^2 + \theta_{el}^2}
\]
and the peak gain is:
\[
  G_{\text{peak}} = \eta \left(\frac{\pi D_{\text{horn}}}{\lambda}\right)^2
\]

Angles are saturated to $\pm 90°$ to avoid gain going negative.

\subsection{Mathematical Analysis}

\subsubsection{Pattern Shape Comparison}

The two models produce fundamentally different pattern shapes:

\begin{table}[H]
\centering
\caption{Horn antenna pattern characteristics.}
\small
\begin{tabularx}{\textwidth}{l X X}
\toprule
\textbf{Property} & \textbf{Pyramidal (\AF{})} & \textbf{cos($\theta$) (\OD{})} \\
\midrule
E/H asymmetry & Yes (independent planes) & No (rotationally symmetric) \\
Sidelobe structure & Nulls and sidelobes from sinc, cos-taper & None (monotonic decay) \\
First E-plane null & $u = \pi$ & Does not exist \\
First H-plane null & $v = 3\pi/2$ & Does not exist \\
Pattern at $\theta = 45°$ & Product of E/H cuts & $\cos(45°) = 0.707$ ($-1.5$~dB) \\
Pattern at $\theta = 89°$ & Near zero (nulls dominate) & $\cos(89°) = 0.017$ ($-17.6$~dB) \\
Physical basis & Aperture diffraction theory & Empirical first-order estimate \\
\bottomrule
\end{tabularx}
\end{table}

\subsubsection{Error Analysis}

The $\cos(\theta)$ model is a reasonable first-order bound for the main beam
envelope but fails to capture:
\begin{enumerate}[leftmargin=2em]
  \item E-plane vs H-plane beamwidth difference (typically 10--30\%)
  \item Sidelobe nulls and peaks (critical for interference calculations)
  \item The broader H-plane beam due to cosine aperture taper
\end{enumerate}

For an X-band standard-gain horn ($G_{\text{peak}} \approx 20$~dBi,
$\theta_{3\text{dB}} \approx 12°$), the $\cos\theta$ model over-estimates gain
at $30°$ off-axis by approximately 8--12~dB compared to the pyramidal model.

\caveat{The pyramidal model captures E/H-plane asymmetry, sidelobe structure,
and null positions, but requires distinct azimuth and elevation beamwidths.
The $\cos(\theta)$ model is simpler and computationally cheaper, producing a
smooth monotonic rolloff suitable for envelope-level simulation but without
sidelobe detail.}

% ══════════════════════════════════════════════════════════════════
\section{Dipole and Monopole Antenna}
\label{sec:dipole}
% ══════════════════════════════════════════════════════════════════

\subsection{Physical Background}

The thin-wire dipole of total length $L$ (in wavelengths) has the classical
far-field pattern:
\begin{equation}
  F(\theta) = \frac{\cos\!\left(\frac{kL}{2}\cos\theta\right)
              - \cos\!\left(\frac{kL}{2}\right)}{\sin\theta}
  \label{eq:dipole_classical}
\end{equation}
where $k = 2\pi/\lambda$ and $\theta$ is measured from the dipole axis.

For a quarter-wave monopole ($L = \lambda/4$) over a perfect ground plane, the
gain is doubled (3~dB increase) in the upper hemisphere, with zero radiation
below the ground plane.

\subsection{\AF{} Implementation}

\AF{} implements the general dipole pattern in \texttt{monopole.py}:
\begin{equation}
  G(\theta_{el}) = G_{\text{peak}} \cdot
    \left[
      \frac{\cos\!\left(\frac{kL}{2}\sin\theta_{el}\right)
            - \cos\!\left(\frac{kL}{2}\right)}
           {\left(1 - \cos\!\left(\frac{kL}{2}\right)\right) \cos\theta_{el}}
    \right]^2
  \label{eq:af_dipole}
\end{equation}

where $\theta_{el}$ is the elevation angle (measured from the horizontal plane,
not the dipole axis).  The denominator includes a normalisation factor
$(1 - \cos(kL/2))$ that ensures the pattern peak is unity, allowing independent
control of peak gain.

\textbf{Edge-case handling:}
\begin{itemize}[leftmargin=2em]
  \item \textbf{Horizon ($|\cos\theta_{el}| < 10^{-10}$):} Returns
        $G_\text{peak} \times 10^{-7}$ to avoid division by zero.
  \item \textbf{Full-wave degenerate ($|1 - \cos(kL/2)| < 10^{-12}$):}
        Falls back to $\cos^n(\theta_{el})$ model, preventing the $0/0$
        singularity that occurs at $L = 2\lambda, 4\lambda, \ldots$ where the
        normalisation factor vanishes.
  \item \textbf{Minimum gain floor:} Output is clamped to
        $G_\text{peak} \times 10^{-7}$ ($-70$~dB relative) to avoid
        numerical underflow.
\end{itemize}

\subsection{\OD{} Implementation}

\OD{} implements the same classical formula separately for dipoles
(\texttt{antenna\_comm\_dipole.cpp}) and monopoles
(\texttt{antenna\_comm\_monopole.cpp}):

\begin{equation}
  G(\theta) = \left[
    \frac{\cos\!\left(\frac{kL}{2}\cos\theta'\right)
          - \cos\!\left(\frac{kL}{2}\right)}{\sin\theta'}
  \right]^2
  \label{eq:od_dipole}
\end{equation}
where $\theta' = \theta_{\text{bse}} - \pi/2$, rotating from the boresight
convention (boresight = along dipole axis) to the formula convention (0 = broadside).

For the monopole, the formula is identical but multiplied by 2 (the ground-plane
gain doubling), with a hard cut to zero for $|\theta_{\text{bse}}| > \pi$:
\begin{equation}
  G_{\text{monopole}}(\theta) = 2 \cdot G_{\text{dipole}}(\theta),
  \qquad |\theta_{\text{bse}}| \leq \pi
\end{equation}

\subsection{Mathematical Analysis}

\subsubsection{Formula Equivalence}

Both implementations use the same classical thin-wire dipole radiation formula.
The apparent difference is purely in \emph{angle convention}:

\begin{table}[H]
\centering
\caption{Dipole angle conventions.}
\small
\begin{tabularx}{\textwidth}{l X X}
\toprule
 & \textbf{\AF{}} & \textbf{\OD{}} \\
\midrule
Input angle & $\theta_{el}$ (from horizon) & $\theta_\text{bse}$ (from dipole axis) \\
Formula angle & $\sin\theta_{el}$ replaces $\cos\theta$ & $\cos(\theta_\text{bse} - \pi/2)$ \\
Identity used & $\cos(\pi/2 - x) = \sin(x)$ & same, applied differently \\
\bottomrule
\end{tabularx}
\end{table}

Setting $\theta_{el} = \pi/2 - \theta$, the two formulations are mathematically
identical.

\subsubsection{Normalisation}

\AF{} normalises by $(1 - \cos(kL/2))$ to decouple the pattern shape from the
peak gain value, enabling the user to specify peak gain independently.  \OD{}
returns the raw un-normalised pattern, where the peak magnitude depends on the
electrical length $L$.  This means \OD{}'s peak gain is implicit in the formula
rather than separately controllable.

\subsubsection{Robustness at Degenerate Points}

At $L = 2\lambda$ (full-wave), $\cos(kL/2) = \cos(2\pi) = 1$, making the
normalisation factor $1 - 1 = 0$.  \AF{} explicitly guards against this with a
$\cos^n$ fallback.  \OD{} does not guard this case; the pattern goes to
$0/0$ as $\theta \to 0$ (boresight) for a full-wave dipole.

Additionally, at $\theta = \pm 90°$ (along the dipole axis), both codes
produce a $0/0$ form.  \AF{} handles this with a magnitude clamp; \OD{} has
an explicit check:
\lstinline[language=C++]{if (abs(abs(angle_off_boresight.value()) - PI/2) < 1e-12) return 0.0;}

\subsubsection{Directive Dipole (\OD{} Only)}

\OD{} includes a \texttt{DirectiveDipole} model that combines the standard
dipole elevation pattern with an azimuthal sinc directivity:
\begin{equation}
  G(\theta_{az}, \theta_{el}) = F_{\text{dipole}}^2(\theta_{el}) \cdot
    G_{\text{az,max}} \cdot \sinc^2\!\left(\theta_{az} \cdot
    \frac{\pi w}{\lambda}\right)
\end{equation}
where $w = 1.22\lambda / \theta_{az,3\text{dB}}$.  This hybrid model provides
azimuthal directivity to what is otherwise an omnidirectional dipole pattern.
\AF{} does not have a direct equivalent, though a similar effect can be achieved
by using the Panel antenna type with a dipole-like elevation beamwidth.

\caveat{The core dipole math is identical between the two implementations.
\AF{} includes explicit guards for full-wave degenerate and horizon
singularity cases.  \OD{} offers a directive dipole variant not present
in \AF{}.  Each approach reflects different design priorities.}

% ══════════════════════════════════════════════════════════════════
\section{Phased Array Antenna}
\label{sec:array}
% ══════════════════════════════════════════════════════════════════

\subsection{Physical Background}

The far-field pattern of a phased array is described by the \emph{Pattern
Multiplication Theorem}:
\begin{equation}
  G_{\text{total}}(\theta, \phi) = G_{\text{element}}(\theta, \phi) \cdot
    |AF(\theta, \phi)|^2
  \label{eq:pmt}
\end{equation}
where the Array Factor for an $N$-element linear array along the $x$-axis is:
\begin{equation}
  AF(\theta) = \sum_{n=0}^{N-1} w_n \, e^{j n (k d \sin\theta + \beta)}
  \label{eq:af_linear}
\end{equation}
with $d$ the element spacing, $\beta$ the progressive phase shift for beam
steering, and $w_n$ the element amplitude weights.

\subsection{\AF{} Implementation}

\AF{} implements the full array factor summation from first principles in
\texttt{array\_antenna.py}:

\begin{equation}
  |AF|^2 = \frac{1}{(\sum w_n)^2}
    \left|\sum_{n=0}^{N-1} w_n \, e^{j(n\psi + \phi_n)}\right|^2,
  \qquad
  \psi = k d \sin\theta + \beta
\end{equation}

\textbf{Supported configurations:}
\begin{itemize}[leftmargin=2em]
  \item \textbf{Geometries:} Linear, planar (rectangular), circular.
  \item \textbf{Planar:} Product of two linear AFs:
    $AF_{\text{planar}} = AF_x(\theta_{az}) \cdot AF_y(\theta_{el})$.
  \item \textbf{Circular:} Elements at $\phi_n = 2\pi n / N$ on radius $R$:
    \[
      AF = \sum_{n=0}^{N-1} w_n \,
        \exp\!\bigl[j k R (\cos(\theta_{az} - \phi_n)\cos\theta_{el}
        - \cos(\theta_{\text{steer}} - \phi_n))\bigr]
    \]
  \item \textbf{Weighting:} Uniform or Taylor (raised-cosine taper for sidelobe
    control).
  \item \textbf{Beam steering:} Progressive phase
    $\beta = -k d \sin\theta_{\text{steer}}$.
  \item \textbf{Element pattern:} Gaussian elliptical beam
    $G_{\text{elem}} = 0.5^{r^2}$.
  \item \textbf{Mutual coupling:} Exponential-decay impedance matrix
    $Z_{mn} = c_0 \, e^{-\alpha |m-n| d/\lambda}$ with Gaussian elimination
    solver for coupled weights $w_{\text{coupled}} = Z^{-1} w_{\text{ideal}}$.
\end{itemize}

\subsection{\OD{} Implementation}

\OD{} models the phased array as a parametric sinc-based rectangular aperture
in \texttt{antenna\_simple\_phased\_array.cpp}:

\begin{equation}
  G = \frac{G_{\text{peak}}}{N_{\text{beams}}} \cdot
    \sinc\!\left((\theta_{az} - \theta_{\text{steer}}) \cdot
    \frac{\pi W}{\lambda}\right) \cdot
    \sinc\!\left((\theta_{el} - \theta_{\text{steer,el}}) \cdot
    \frac{\pi H}{\lambda}\right)
  \label{eq:od_array}
\end{equation}

\textbf{Key features:}
\begin{itemize}[leftmargin=2em]
  \item \textbf{Multi-beam:} Supports TX, RX, and TX/RX (TRX) beams with
    independent steering angles. TX gain divided by number of TX beams.
  \item \textbf{Spoiled transmit:} Per-beam configurable effective aperture
    dimensions for broadened transmit coverage.
  \item \textbf{Beam management:} Add, resteer, and remove beams at runtime
    (essential for system simulation).
  \item \textbf{Scan loss:} Optional $\cos(\theta_{\text{steer}})$ scan loss
    when enabled.
\end{itemize}

\subsection{Mathematical Analysis}

\subsubsection{Pattern Fidelity}

\begin{table}[H]
\centering
\caption{Phased array model comparison.}
\small
\begin{tabularx}{\textwidth}{l X X}
\toprule
\textbf{Property} & \textbf{AF Phasor Sum (\AF{})} & \textbf{Sinc Aperture (\OD{})} \\
\midrule
Main beam shape & Exact for discrete elements & Continuous-aperture approximation \\
Grating lobes & Naturally emerge at $d > \lambda/2$ & Not captured \\
Sidelobe structure & Correct for $N$ elements + taper & Fixed sinc sidelobes \\
Taylor taper effect & Reduced sidelobes, wider main beam & Not available \\
Mutual coupling & Perturbed amplitude and phase & Not modelled \\
Scan loss & Via $\cos\theta$ element factor & Optional toggle \\
Computation cost & $O(N)$ per angle & $O(1)$ per angle \\
\bottomrule
\end{tabularx}
\end{table}

The critical difference: for a discrete array with $d/\lambda > 0.5$, grating
lobes appear at angles where $kd\sin\theta = 2\pi m$.  The AF phasor sum
naturally produces these; the sinc aperture model cannot.  This can lead to
significant under-estimation of interference in system simulations when
element spacing exceeds half a wavelength.

\subsubsection{Multi-Beam Operation}

\OD{}'s multi-beam management is a significant capability for system-level
simulation.  The total transmitted power is split equally among active TX beams
($G_{\text{tx}}/N_{\text{beams}}$), and modern digital beamforming is assumed
for receive (no RX beam-split loss).  \AF{} does not currently model
multi-beam operation.

\subsubsection{Circular Array}

\AF{}'s circular array uses the geometrically exact phase delay for each element
on the ring, correctly capturing the non-uniform angular spacing of grating
lobe-like artefacts characteristic of circular arrays.  \OD{} does not offer a
circular geometry.

\caveat{The first-principles array factor naturally produces grating lobes,
taper effects, and coupling perturbations at $O(N)$ cost per point.  The
parametric sinc model runs at $O(1)$ per point and supports multi-beam
management, but does not capture grating lobes or discrete-element effects.
The trade-off is pattern detail vs.\ computational speed and system-level
features.}

% ══════════════════════════════════════════════════════════════════
\section{Panel / Sector Antenna}
\label{sec:panel}
% ══════════════════════════════════════════════════════════════════

\subsection{Physical Background}

A rectangular aperture of width $a$ and height $b$ with uniform illumination
produces a pattern that is the product of two sinc functions:
\begin{equation}
  G(\theta_{az}, \theta_{el}) = G_\text{peak} \cdot
    \sinc^2\!\left(\frac{\pi a}{\lambda}\sin\theta_{az}\right) \cdot
    \sinc^2\!\left(\frac{\pi b}{\lambda}\sin\theta_{el}\right)
\end{equation}

\subsection{\AF{} Implementation}

\AF{} implements this model in \texttt{panel.py} with beamwidth-normalised
arguments:
\begin{equation}
  G = G_\text{peak} \cdot \sinc^2(u_{az}) \cdot \sinc^2(u_{el}) \cdot
    f_\text{front}
\end{equation}
where
\[
  u_{az} = 1.3916 \cdot \frac{\sin\theta_{az}}{\sin(\theta_{az,3\text{dB}}/2)},
  \qquad
  u_{el} = 1.3916 \cdot
    \frac{\sin(\theta_{el} - \theta_\text{tilt})}
         {\sin(\theta_{el,3\text{dB}}/2)}
\]

The constant $K = 1.3916$ satisfies $\sinc^2(K) = 0.5$, ensuring exact $-3$~dB
at the declared beamwidth.

Additional features:
\begin{itemize}[leftmargin=2em]
  \item \textbf{Electrical downtilt:} Configurable beam tilt in elevation.
  \item \textbf{FTB modelling:} Same continuous back-lobe model as all
    directional types.
\end{itemize}

\subsection{Comparison with \OD{}}

\OD{} does not have a dedicated panel/sector antenna type.  However, the
\texttt{ParabolicEllipticalDish} uses a separable sinc model with independent
width and height:
\begin{equation}
  G = \eta \cdot \frac{4\pi A}{\lambda^2} \cdot
    |\sinc(\theta_{az} \cdot \pi W f / c)| \cdot
    |\sinc(\theta_{el} \cdot \pi H f / c)|
\end{equation}
where $A = \pi W H / 4$ (elliptical area).  This is mathematically similar to
\AF{}'s Panel model but assumes an elliptical rather than rectangular aperture
and lacks the beamwidth normalisation that guarantees the $-3$~dB points align
with declared beamwidth.

\caveat{\OD{}'s elliptical dish uses $|\sinc|$ (magnitude, not squared),
which means its sidelobe envelope decays as $1/u$ rather than $1/u^2$ --- the
sidelobes are significantly higher than either the squared-sinc model (\AF{}) or
physical measurement.}

% ══════════════════════════════════════════════════════════════════
\section{LPDA Antenna}
\label{sec:lpda}
% ══════════════════════════════════════════════════════════════════

\subsection{\AF{} Implementation}

The Log-Periodic Dipole Array is modelled in \texttt{lpda.py} using a Gaussian
elliptical main beam with radial sidelobe envelope:

\begin{equation}
  G_\text{main} = G_\text{peak} \cdot 0.5^{r^2} \cdot f_\text{front},
  \qquad
  r^2 = \left(\frac{\theta_{az}}{\theta_{az,3\text{dB}}/2}\right)^2 +
        \left(\frac{\theta_{el}}{\theta_{el,3\text{dB}}/2}\right)^2
\end{equation}

The Gaussian beam model $0.5^{r^2}$ is equivalent to:
\[
  G = G_\text{peak} \cdot \exp(-\ln 2 \cdot r^2)
\]
which produces smooth elliptical iso-gain contours.  This is a pragmatic
engineering choice: LPDAs do not have a simple closed-form pattern, so the
Gaussian model matches the main beam shape well while the sidelobe envelope
provides the correct far-out behaviour.

\textbf{Sidelobe model:}
\begin{itemize}[leftmargin=2em]
  \item Radial (elliptical rings, not rectangular bands).
  \item Configurable first sidelobe level, number, and decay rate.
  \item Smooth taper beyond last sidelobe (no hard cutoff).
  \item Envelope: $G_\text{sl} = A_1 \cdot d^n \cdot \cos^2(\text{phase})$
    with exponential decay beyond last lobe.
\end{itemize}

\OD{} does not implement an LPDA antenna type.

% ══════════════════════════════════════════════════════════════════
\section{Omnidirectional / Isotropic Antenna}
\label{sec:omni}
% ══════════════════════════════════════════════════════════════════

\subsection{\AF{} Omni Implementation}

\AF{}'s omnidirectional antenna uses a $\cos^n$ elevation model with sidelobe
envelope:
\begin{equation}
  G(\theta_{el}) = G_\text{peak} \cdot |\cos\theta_{el}|^{n_{el}},
  \qquad
  n_{el} = \frac{\ln 0.5}{\ln \cos(\theta_{el,3\text{dB}}/2)}
\end{equation}

This models the elevation shaping of real-world omnidirectional antennas
(collinear arrays, ground-plane monopoles, discone antennas) while maintaining
360° azimuth coverage.

\subsection{\OD{} Isotropic Implementation}

\OD{}'s isotropic antenna returns $G = 1$ at all angles, serving as a
theoretical reference.

\subsection{Comparison}

These serve different purposes: \AF{}'s Omni models a \emph{real-world}
omnidirectional antenna with elevation shaping, while \OD{}'s Isotropic is
a \emph{theoretical} point source.  For an ideal isotropic reference in \AF{},
setting $n_{el} = 0$ and $G_\text{peak} = 1$ gives the equivalent result.

% ══════════════════════════════════════════════════════════════════
\section{System-Level Features}
\label{sec:features}
% ══════════════════════════════════════════════════════════════════

Beyond the core antenna patterns, the two frameworks differ in their
system-level modelling capabilities.  These features are summarised below.

\subsection{Features Unique to \AF{}}

\begin{description}[leftmargin=1.5em, style=nextline]
  \item[Front-to-Back Ratio] Every directional antenna type includes a
    configurable FTB with continuous back-lobe weighting
    $(1 - \cos\theta_{az})/2$.  This ensures no dead zone at $\pm 90°$
    and physically plausible rear-hemisphere radiation.

  \item[Sigmoid Front Fade] The transition between front and back hemispheres
    uses a logistic sigmoid $f = 1/(1 + e^{-k\cos\theta})$ whose steepness
    $k$ is user-configurable (default~$40$, range~$1$--$100$).  At the
    default $k = 40$ the transition zone spans $\sim 10°$; reducing $k$ to
    $\sim 12$ widens it to $\sim 30°$, better matching measured horn and
    dish roll-off.  This eliminates the visible discontinuity at $\pm 90°$
    in heatmap plots.

  \item[Frequency-Dependent Beamwidth] Three modes:
    \begin{enumerate}
      \item Exponent: $\theta_\text{BW}(f) = \theta_\text{BW,ref} \cdot (f_\text{ref}/f)^\alpha$
      \item Percentage: $X\%$ decay per octave from reference
      \item Three-point: Quadratic Lagrange through $(f_\text{min}, f_\text{mid}, f_\text{max})$
    \end{enumerate}

  \item[VSWR Rolloff] Band-edge gain reduction with configurable depth,
    bandwidth fraction, and shape (linear or cosine).

  \item[Pattern Breakup] Deterministic multi-frequency ripple at high off-axis
    angles, modelling diffraction and structural scattering.

  \item[Ground Reflection] Two-ray interference model:
    $G' = G \cdot (1 + \rho^2 + 2\rho\cos(\Delta\phi + \pi))$.

  \item[Cross-Polarisation] Angular-dependent cross-pol with configurable
    boresight isolation, peak angle, and diagonal-plane factor.

  \item[Asymmetry] Squint (az) and tilt (el) offsets with optional
    deterministic perturbation for realistic imperfections.

  \item[Sidelobe Envelope] Configurable first sidelobe level, number of
    sidelobes, decay rate (Taylor vs uniform), and radial (elliptical)
    geometry.

  \item[Mutual Coupling (Array)] Impedance matrix model with Gaussian
    elimination solver perturbing element amplitudes and phases.

  \item[Gain-Beamwidth Coupling] Frequency-dependent gain linked to
    beamwidth via directivity relationships, with three configurable
    modes (gain drives beamwidth, beamwidth drives gain, independent
    curves).
\end{description}

\subsection{Features Unique to \OD{}}

\begin{description}[leftmargin=1.5em, style=nextline]
  \item[Monopulse Receive] Sum ($\Sigma$) and difference ($\Delta_{az}$,
    $\Delta_{el}$) channels on dish and elliptical dish antennas, using
    standard monopulse comparator approximations.

  \item[Multi-Beam Phased Array] Runtime beam management (add, resteer, remove)
    with TX/RX/TRX beam types, spoiled transmit beams, and automatic power
    splitting.

  \item[Lookup-Based Antennas] 2-D (separate az/el cuts) and 3-D (full
    az/el/freq) table-interpolated measured patterns.  Beamwidth is
    determined from the interpolated data by searching for the $-3$~dB points.

  \item[Directive Dipole] Hybrid dipole elevation $\times$ sinc azimuth model.
\end{description}

% ══════════════════════════════════════════════════════════════════
\section{Numerical Robustness and Edge Cases}
\label{sec:robustness}
% ══════════════════════════════════════════════════════════════════

\begin{longtable}{p{3.5cm} p{5.5cm} p{5.5cm}}
\toprule
\textbf{Edge Case} & \textbf{\AF{}} & \textbf{\OD{}} \\
\midrule
\endfirsthead
\toprule
\textbf{Edge Case} & \textbf{\AF{}} & \textbf{\OD{}} \\
\midrule
\endhead
\midrule
\multicolumn{3}{r}{\itshape Continued on next page} \\
\bottomrule
\endfoot
\bottomrule
\endlastfoot

sinc(0)/jinc(0) &
  Explicit clamp: returns 1.0 when $|u| < 10^{-10}$ &
  Uses built-in sinc; $0/0$ can occur at exact boresight \\

Dipole at $\theta = \pm 90°$ &
  Returns $G_\text{peak} \times 10^{-7}$ (floor) &
  Returns 0.0 (hard check $< 10^{-12}$) \\

Full-wave dipole ($kL = 4\pi$) &
  $\cos^n$ fallback when $|1 - \cos(kL/2)| < 10^{-12}$ &
  Not guarded; potential $0/0$ at boresight \\

Horn at $v = \pm\pi/2$ &
  L'H\^{o}pital limit: $\pi/4 \cdot \cos(v)$ &
  N/A (no H-plane function) \\

Negative gain &
  All types clamp to $G_\text{peak} \times 10^{-7}$ minimum &
  No universal floor; some types can return 0 or near-zero \\

Array with $N = 0$ elements &
  Weights list is empty; $|AF|^2 = 0$ &
  Not guarded (likely division by zero) \\

Horn angle saturation &
  Sigmoid fade (configurable $k$) handles gracefully &
  Hard saturation at $\pm 90°$: \texttt{Saturate()} \\
\end{longtable}

% ══════════════════════════════════════════════════════════════════
\section{Performance Considerations}
\label{sec:perf}
% ══════════════════════════════════════════════════════════════════

\begin{table}[H]
\centering
\caption{Computational complexity per pattern point.}
\small
\begin{tabularx}{\textwidth}{l X X}
\toprule
\textbf{Type} & \textbf{\AF{}} & \textbf{\OD{}} \\
\midrule
Dish & $O(1)$: one Bessel + arithmetic & $O(1)$: one sinc + arithmetic \\
Horn & $O(1)$: two trig functions & $O(1)$: one trig function \\
Dipole & $O(1)$: four trig functions & $O(1)$: four trig functions \\
Phased Array & $O(N)$: phasor sum over $N$ elements & $O(1)$: two sinc evaluations \\
Panel & $O(1)$: two sinc evaluations & N/A \\
LPDA & $O(1)$: exp + sidelobe search & N/A \\
\bottomrule
\end{tabularx}
\end{table}

\AF{} compensates for the $O(N)$ array factor cost with NumPy vectorisation,
computing the entire $(az, el)$ grid in batched array operations (typically
50--100$\times$ faster than the scalar Python fallback).  \OD{}'s $O(1)$
parametric model is advantageous for real-time system simulation where thousands
of link calculations per timestep are needed.

% ══════════════════════════════════════════════════════════════════
\section{Conclusions and Recommendations}
\label{sec:conclusions}
% ══════════════════════════════════════════════════════════════════

\subsection{Pattern Modelling Approaches}

The two codebases take different approaches to pattern modelling for each
shared antenna type:

\begin{enumerate}[leftmargin=2em]
  \item \textbf{Dish:} \AF{} uses the Jinc$^2$ (Airy) kernel for circular
    aperture diffraction; \OD{} uses a sinc kernel derived from rectangular
    aperture theory.  These produce different sidelobe levels and decay rates.
  \item \textbf{Horn:} \AF{} models separate E/H-plane aperture distributions
    for a pyramidal horn; \OD{} uses a rotationally symmetric $\cos(\theta)$
    envelope for a conical horn.  These target different horn geometries.
  \item \textbf{Dipole:} Both use the same classical thin-wire formula with
    different angle conventions and normalisation strategies.
  \item \textbf{Array:} \AF{} computes a first-principles phasor sum (with
    grating lobes and taper); \OD{} uses a parametric sinc aperture model
    (with multi-beam management and $O(1)$ performance).
\end{enumerate}

\subsection{System-Level Capability}

\OD{} provides system-simulation features not present in \AF{}:
\begin{itemize}[leftmargin=2em]
  \item Monopulse tracking (sum/difference channels)
  \item Multi-beam phased array with runtime beam management
  \item Measured-pattern lookup tables (2-D and 3-D)
  \item Directive dipole hybrid model
\end{itemize}

\AF{} provides post-processing features not present in \OD{}:
\begin{itemize}[leftmargin=2em]
  \item Front-to-back ratio with continuous back-lobe weighting
  \item Frequency-dependent beamwidth (three configurable modes)
  \item VSWR rolloff, pattern breakup, cross-polarisation
  \item Ground reflection (two-ray model)
  \item Mutual coupling for phased arrays
\end{itemize}

The two feature sets reflect different design priorities: \OD{} targets
real-time system-of-systems simulation, while \AF{} targets detailed
pattern generation and analysis.

\subsection{Potential Enhancements}

Based on this analysis, the following cross-pollination opportunities exist:

\begin{table}[H]
\centering
\caption{Recommended enhancements.}
\small
\begin{tabularx}{\textwidth}{l X X}
\toprule
\textbf{Enhancement} & \textbf{Target} & \textbf{Source} \\
\midrule
Lookup-based antenna & \AF{} & \OD{} (2-D/3-D tables) \\
Monopulse receive mode & \AF{} & \OD{} (dish, array) \\
Multi-beam array & \AF{} & \OD{} (beam management) \\
Directive dipole type & \AF{} & \OD{} (dipole $\times$ sinc) \\
Jinc$^2$ dish model & \OD{} & \AF{} (Airy pattern) \\
Pyramidal horn model & \OD{} & \AF{} (E/H aperture) \\
FTB back-lobe model & \OD{} & \AF{} (continuous weighting) \\
Freq-dependent BW & \OD{} & \AF{} (3 decay modes) \\
Pattern breakup & \OD{} & \AF{} (deterministic ripple) \\
\bottomrule
\end{tabularx}
\end{table}

% ══════════════════════════════════════════════════════════════════
\appendix
\section{Mathematical Reference}
\label{app:math}
% ══════════════════════════════════════════════════════════════════

\subsection{Bessel Function $J_1(x)$}

The first-order Bessel function of the first kind is defined by:
\begin{equation}
  J_1(x) = \frac{x}{2} \sum_{m=0}^{\infty}
    \frac{(-1)^m}{m!\,(m+1)!} \left(\frac{x}{2}\right)^{2m}
    = \frac{x}{2} - \frac{x^3}{16} + \frac{x^5}{384} - \cdots
\end{equation}

\AF{}'s fallback approximation uses the rational polynomial fit from
Abramowitz \& Stegun (1964), \S9.4, providing $\sim$8 significant digits
for $|x| < 8$.

\subsection{Sinc Function Convention}

Both codebases use the \emph{unnormalised} sinc:
\[
  \sinc(x) = \frac{\sin(x)}{x}, \qquad \sinc(0) = 1
\]

Note: some references use the \emph{normalised} sinc $\sin(\pi x)/(\pi x)$.
This report and both codebases use the unnormalised form throughout.

\subsection{$\cos^n$ Beam Model}

The exponent $n$ is determined from the half-power beamwidth $\theta_{3\text{dB}}$:
\[
  \cos^n(\theta_{3\text{dB}}/2) = 0.5
  \quad\Longrightarrow\quad
  n = \frac{\ln 0.5}{\ln \cos(\theta_{3\text{dB}}/2)}
\]

\subsection{Front-Hemisphere Sigmoid}

The sigmoid transition function:
\[
  f(\cos\theta_{az}) = \frac{1}{1 + e^{-k\cos\theta_{az}}}
\]
gives $f = 0.5$ at $\theta_{az} = \pm 90°$ for any $k$.  The steepness
parameter $k$ is user-configurable (default~$40$, range~$1$--$100$) and
controls the width of the front/back transition zone:
\begin{center}
\begin{tabular}{rll}
  $k$ & Transition zone ($3\%$--$97\%$) & Character \\
  \hline
  $10$ & $\sim 40°$ & Gentle, physical \\
  $20$ & $\sim 20°$ & Moderate \\
  $40$ & $\sim 10°$ & Sharp (default) \\
  $80$ & $\sim 5°$  & Very sharp \\
\end{tabular}
\end{center}
Most measured directive antennas exhibit a gradual roll-off through
$\pm 90°$; reducing $k$ to $10$--$15$ often produces more realistic
heatmaps, particularly for wide-beam horns and dishes.

\subsection{Array Factor Normalisation}

For uniform weights $w_n = 1$, the array factor of an $N$-element linear array
simplifies to:
\[
  |AF|^2 = \frac{1}{N^2}
    \left|\frac{\sin(N\psi/2)}{\sin(\psi/2)}\right|^2,
  \qquad
  \psi = kd\sin\theta + \beta
\]
Grating lobes appear when $\psi = 2\pi m$ for integer $m \neq 0$, i.e.\
when $\sin\theta = -\beta/(kd) + m\lambda/d$.

\end{document}
